\documentclass[usenames,dvipsnames]{beamer}

\usepackage{xcolor}
\usepackage[utf8]{inputenc}
\usepackage{listings}
\usepackage{tikz}
\usepackage[style=alphabetic,giveninits=true,terseinits=true]{biblatex}
\usepackage{graphicx}

\usetikzlibrary{matrix, positioning, decorations.pathreplacing, overlay-beamer-styles}

\definecolor{tumblue}{RGB}{0,101,189}
\definecolor{dkgreen}{rgb}{0,0.6,0}
\definecolor{mauve}{rgb}{0.58,0,0.82}
\lstset{
    basicstyle=\footnotesize\ttfamily,%
    columns=fullflexible,%
    keywordstyle=\bfseries\color{blue},%
    commentstyle=\color{dkgreen},%
    stringstyle=\color{mauve},%
}

\DeclareFieldFormat{url}{\href{#1}{\resizebox{!}{2ex}{\lower3pt\hbox{\pgfuseimage{beamericononline}}}}}
\addbibresource{slides.bib}

\setbeamerfont{footnote}{size=\tiny}
\renewcommand\footnoterule{}
\usecolortheme{orchid}
\setbeamertemplate{navigation symbols}{}
\setbeamertemplate{sidebar right}{% also implies no navbar
    \llap{\tikz\fill[tumblue,scale=.1] (0,0) ++(-5cm,-5cm) ++(-10cm,0)
        +(0cm,0cm) -- +(4cm,0cm) -- +(4cm,-4cm) -- +(5cm,-4cm) -- +(5cm,0cm) --
        +(10cm,0cm) -- +(10cm,-5cm) -- +(9cm,-5cm) -- +(9cm,-1cm) --
        +(8cm,-1cm) -- +(8cm,-5cm) -- +(7cm,-5cm) -- +(7cm,-1cm) --
        +(6cm,-1cm) -- +(6cm,-5cm) -- +(3cm,-5cm) -- +(3cm,-1cm) --
        +(2cm,-1cm) -- +(2cm,-5cm) -- +(1cm,-5cm) -- +(1cm,-1cm) --
        +(0cm,-1cm) -- cycle;}
    % Slide number in triangle.
    \vfill\llap{\tikz\path[fill=tumblue,inner sep=3pt] (0,0) -- ++(-1cm,0) -- +(1cm,1cm) -- cycle node[anchor=south east] {\usebeamerfont{footline}\color{white}\bfseries\insertframenumber};}%
% Variant without triangle
%\vfill\llap{\usebeamerfont{footline}\usebeamercolor[fg]{footline}\insertframenumber\hskip3pt}\vskip3pt%
}

\graphicspath{ {./images/} }

\title{CoPaaS - Compiler-oriented Programming as a Service}
\author{Simon Bußmann \and Johannes Maier}
\institute{School of Computation, Information, and Technology\\Technical University of Munich}
\date{Future}

\begin{document}

    \begin{frame}
        \titlepage
    \end{frame}

    \begin{frame}{The program}
        \begin{itemize}
            \item JIT-Compiler for toy ISA
            \item Instructions: \lstinline{ADD, ADDI, SUB, COPY, LOADI}
            \item Important: No jumps, branches, function calls, etc. for simplicity
            \item Registers:
            \only<1> {
                \begin{tabular}{|l|l|l|}
                    \hline
                    \textbf Adelheid  & \textbf Friedrich & \textbf Klaus    \\
                    \hline
                    \textbf Berthold  & \textbf Giesela   & \textbf Liesbeth \\
                    \hline
                    \textbf Cornelia  & \textbf Heinrich  & \textbf Manfred  \\
                    \hline
                    \textbf Dora      & \textbf Irmgard   & \textbf Norbert  \\
                    \hline
                    \textbf Engelbert & \textbf Joachim   &                  \\
                    \hline
                \end{tabular}
            }
            \only<2-> {
                \texttt A, \texttt B\dots \texttt N
            }
            \visible<3-> {
                \item Seccomp jail for generated code
                \item Except: \visible<4->{In premium mode activated by key}
            }
            \visible<5->{
                \item NX, PIE, Full RERLO, Stack Canary
            }
        \end{itemize}
    \end{frame}

    \begin{frame}{x86-64 Instruction Encoding Basics}
        \begin{columns}
            \begin{column}{0.35\textwidth}

                Parts of an instruction:
                \begin{itemize}
                    \item \textcolor<2>{blue}{Prefixes}
                    \item Opcode
                    \item \textcolor<3-5>{blue}{ModR/M}
                    \item SIB
                    \item Displacement
                    \item Immediate
                \end{itemize}
            \end{column}
            \begin{column}{0.65\textwidth}
                \begin{tikzpicture}[scale=0.75]

                    % REX.W prefix
                    \matrix (rex) at (0,0) [
                    matrix of nodes,
                    nodes={draw, minimum size=18pt, anchor=center},
                    column sep=-\pgflinewidth,
                    row sep=-\pgflinewidth,
                    ampersand replacement=\&
                    ] {
                        0 \& 1\& 0 \& 0 \& \textcolor<2>{blue}{W} \& \textcolor<2>{blue}{R} \& X \& \textcolor<2>{blue}{B} \\
                    };

                    % ModR/M byte
                    \matrix (modrm) [
                        below=20pt of rex,
                        matrix of nodes,
                        nodes={draw, minimum size=18pt, anchor=center},
                        column sep=-\pgflinewidth,
                        row sep=-\pgflinewidth,
                        ampersand replacement=\&
                    ] {
                        \textcolor<3>{blue}{1} \& \textcolor<3>{blue}{1} \& \textcolor<4>{blue}{0} \& \textcolor<4>{blue}{1} \& \textcolor<4>{blue}{0} \& \textcolor<5>{blue}{0} \& \textcolor<5>{blue}{0} \& \textcolor<5>{blue}{0} \\
                    };

                    % Labels
                    \node[anchor=east] at (rex.west) {REX.W};
                    \node[anchor=east] at (modrm.west) {ModR/M};

                    % Brackets for ModR/M parts
                    \draw[semithick, decorate, decoration={brace, amplitude=10pt, mirror, raise=4pt}, alt=<3>{blue}{black}] (modrm-1-1.south west) -- (modrm-1-2.south east) node[midway, below=14pt] {\textcolor<3>{blue}{Modifier}};
                    \draw[semithick, decorate, decoration={brace, amplitude=10pt, mirror, raise=4pt}, alt=<4>{blue}{black}] (modrm-1-3.south west) -- (modrm-1-5.south east) node[midway, below=14pt, xshift=-4pt] {\textcolor<4>{blue}{Reg}};
                    \draw[semithick, decorate, decoration={brace, amplitude=10pt, mirror, raise=4pt}, alt=<5>{blue}{black}] (modrm-1-6.south west) -- (modrm-1-8.south east) node[midway, below=14pt] {\textcolor<5>{blue}{RM}};

                \end{tikzpicture}
            \end{column}
        \end{columns}
    \end{frame}

    \begin{frame}{The Bug}
        \only<2-3> {
            \begin{figure}
                \centering
                \includegraphics[scale=0.65]{bug}
            \end{figure}
        }
        \visible<3>{No, not you...}
        \only<4->{
            \begin{itemize}
                \item Faulty encoding of instructions using \lstinline{r[8-15]}
                \item[$\Rightarrow$] Register ids are 4bit long
                \item Example: encode ModR/M for \lstinline{mov r8, r9}
                \item Register ids: $\text r8 = 0\text b1000, \text r9 = 0\text b1001$
                \[\text{Mod R/M} = \textcolor<4>{blue}{0\text{b}1100\_0000} + \textcolor<5>{blue}{(\text{reg2\_id} \ll 3)} + \textcolor<6>{blue}{\text{reg1\_id}}\]
            \end{itemize}

            \begin{center}
                \begin{tikzpicture}
                    \matrix (modrm) [
                        below=20pt of rex,
                        matrix of nodes,
                        nodes={draw, minimum size=18pt, anchor=center},
                        column sep=-\pgflinewidth,
                        row sep=-\pgflinewidth,
                        ampersand replacement=\&
                    ] {
                        \textcolor<4>{blue}{\textcolor<5>{red}{\alt<5->{0}{1}}} \&
                        \textcolor<4>{blue}{\textcolor<5>{red}{\alt<5->{0}{1}}} \&
                        \textcolor<5>{blue}{0} \&
                        \textcolor<5>{blue}{\textcolor<6>{red}{\alt<6->{1}{0}}} \&
                        \textcolor<5>{blue}{\textcolor<6>{red}{\alt<5>{1}{0}}} \&
                        \textcolor<6>{blue}{0} \&
                        \textcolor<6>{blue}{0} \&
                        \textcolor<6>{blue}{0} \\
                    };
                \end{tikzpicture}
            \end{center}
        }
        \only<6>{
            \begin{itemize}
                \item Actually encoded instruction: \lstinline{mov [r8], r10}
            \end{itemize}
        }
    \end{frame}

    \begin{frame}{Exploit Primitives}
        But, what can we do with that?
        % TODO: describe primitives further
        \begin{itemize}
            \item Arbitrary write: \lstinline{mov [reg1], reg2}
            \only<1-2> {
                \begin{itemize}
                    \item Use \lstinline{COPY} to generate \lstinline{mov}, overflow into modifier bits
                    \item Register id constraints: $\text{reg1\_id} = 0\text{b}0\text{xxx}, \text{reg2\_id}=0\text{b}1\text{xxx}$
                    \item[$\Rightarrow$] But can rename registers freely
                    \item Example: \texttt{mov [\textcolor{ForestGreen}{rax}], \textcolor{blue}{r8}}
                \end{itemize}

                \begin{center}
                    \begin{tikzpicture}
                        \matrix (modrm) [
                            below=20pt of rex,
                            matrix of nodes,
                            nodes={draw, minimum size=18pt, anchor=center},
                            column sep=-\pgflinewidth,
                            row sep=-\pgflinewidth,
                            ampersand replacement=\&
                            ] {
                            \textcolor<1>{black}{\textcolor<2>{red}{\alt<2>{0}{1}}} \&
                            \textcolor<1>{black}{\textcolor<2>{red}{\alt<2>{0}{1}}} \&
                            \textcolor<2>{blue}{0} \&
                            \textcolor<2>{blue}{0} \&
                            \textcolor<2>{blue}{0} \&
                            \textcolor<2>{ForestGreen}{0} \&
                            \textcolor<2>{ForestGreen}{0} \&
                            \textcolor<2>{ForestGreen}{0} \\
                            };
                    \end{tikzpicture}
                \end{center}
            }

            \item Arbitrary read: \lstinline{mov reg1, [reg2]}
            \only<3>{
                \begin{itemize}
                \item Actually: \lstinline{mov reg1, 0; add reg1, [reg2]}
                \item Use same trick as for write
                \item There exists multiple opcodes for add\footnote{Also for most other standard x86 operations.}, for example:
                \item Opcode $= 0x1$: \lstinline{mov [reg1], reg2}
                \item Opcode $= 0x3$: \lstinline{mov reg1, [reg2]}
                \end{itemize}
            }

            \item Stack access
            \only<4> {
                \begin{itemize}
                \item Another overflow, but this time into \lstinline{reg2_id}
                \item Can create \lstinline{reg2_id + 1}, using reg1\_id = $0\text{b}1$xxx
                \item Id of \lstinline{rsp} is $0\text{b}0100$
                \item[$\Rightarrow$] For \lstinline{COPY G, B} is actually \lstinline{mov r8, rsp} generated
                \item Combined allow those primitives unlimited access to the stack
                \end{itemize}
            }
        \end{itemize}
    \end{frame}

    \begin{frame}{The Exploit}
        \only<1-4>{
        Steps:
        \begin{itemize}
            \item Leak activation key
            \only<2>{
                \begin{itemize}
                    \item Activation key is stored in a variable in the binary
                    \item The provided program returns into the binary after execution
                    \item Use stack access primitive to leak the return address
                    \item Use arbitrary read primitive to leak the key byte by byte
                \end{itemize}
            }
            \item Escape seccomp jail
            \only<3>{
                \begin{itemize}
                    \item Activate premium mode with the leaked key 
                    \item \it{much wow great hacking}
                \end{itemize}
            }
            \item Execute \lstinline{system("/bin/sh")}
            \only<4>{
                \begin{itemize}
                    \item Use stack access primitive to leak a libc address
                    \item Calculate \lstinline{system} and \lstinline{/bin/sh} addresses
                    \item Load \lstinline{/bin/sh} address into \lstinline{rdi}
                    \item Use arbitrary write primitive to overwrite saved rip
                \end{itemize}
            }
        \end{itemize}
        }
        \only<5>{
            \begin{figure}
                \centering
                \includegraphics[scale=0.22]{profit_penguin}
            \end{figure}
        }
    \end{frame}

\end{document}
